\begin{abstract}
Emerging studies, exemplified by BitNet~\cite{bitnet}, are ushering in a transformative phase for Large Language Models (LLMs) based on 1-bit representations. In this paper, we present \textbf{\text{BitNet b1.58}}, a ternary-parameter LLM wherein each individual weight can assume values from the set \{-1, 0, 1\}. This architecture achieves parity with standard full-precision Transformers (FP16 or BF16) of identical scale and trained on equivalent data, matching them in perplexity and downstream task effectiveness, yet offering substantially improved latency, memory efficiency, throughput, and reduced energy demands. Notably, the proposed 1.58-bit LLM introduces a novel scaling law and training methodology for crafting LLMs that are simultaneously resource-efficient and high-performing. This innovation also supports a new computational paradigm and paves the way for the development of hardware architectures specifically optimized for 1-bit LLM execution.
\end{abstract}

\section{The Era of 1-bit LLMs}

The advancement of artificial intelligence has recently been characterized by the development and proliferation of ever-larger Large Language Models (LLMs), whose heightened capabilities have driven significant progress across diverse natural language processing domains. Nevertheless, this expansion in model size introduces substantial hurdles for deployment and intensifies both economic and environmental concerns, particularly in relation to energy usage.

To confront these issues, a common strategy involves the utilization of post-training quantization to derive low-bitwidth models for inference~\cite{smoothquant, gptq, quip, quip_sharp}. By lowering the numerical precision of both weights and activations, this approach considerably lessens the computational burden and memory footprint of LLMs. Industry practices increasingly favor shifting from 16-bit formats towards lower-precision solutions, such as 4-bit configurations~\cite{gptq, awq}. However, despite its prevalence in commercial LLM deployments, post-training quantization exhibits limitations and generally falls short of optimality.

Recent advances in 1-bit neural network designs, exemplified by BitNet~\cite{bitnet}, point toward a compelling strategy for cutting down the operational costs of large language models (LLMs) while largely preserving their effectiveness. Standard LLMs operate using 16-bit floating-point numbers (such as FP16 or BF16), and the dominant computational workload in these models stems from matrix multiplications, which primarily rely on floating-point addition and multiplication. BitNet diverges from this by restricting matrix multiplications to integer addition only, leading to substantial reductions in energy consumption. Given that power is often the limiting factor for computational throughput on many hardware platforms, these energy savings naturally yield speed improvements as well.

Apart from the computation itself, significant overhead arises during inference from shuttling model parameters between DRAM and faster, but more costly, on-chip memory such as SRAM. While scaling up SRAM capacity can boost throughput, it typically results in steep cost increases relative to DRAM. Compared to their full-precision counterparts, 1-bit LLMs demand far less in terms of both memory storage and bandwidth, making parameter loading more efficient and thereby facilitating quicker, less expensive inference.

In this paper, we present \textbf{\text{BitNet b1.58}}, an important extension of the 1-bit LLM framework. In this variant, each network parameter can take on one of three possible values: \{-1, 0, 1\}, thus introducing a zero in addition to the two values used in the original 1-bit BitNet and raising the effective bit-width to 1.58 in binary encoding. \text{BitNet b1.58} preserves all advantages of the original BitNet architecture, such as its unique computational approach—nearly eliminating the need for multiplication during matrix multiplications—and its high efficiency in energy usage, memory footprint, throughput, and inference latency when set against FP16 LLM standards. Beyond these retained properties, \text{BitNet b1.58} brings two notable enhancements. First, it strengthens model expressiveness by allowing explicit feature suppression through zero-valued weights, leading to improved outcomes in low-bit LLM scenarios. Second, our empirical results indicate that \text{BitNet b1.58} is capable of achieving parity with full-precision (FP16) models in both perplexity and downstream task performance from scales of 3B parameters onward, given identical settings such as model size and number of training tokens.

\section{BitNet b1.58}

\text{BitNet b1.58} builds upon the BitNet framework, a Transformer variant where the conventional \emph{nn.Linear} layers are substituted by \emph{BitLinear} modules. The model is trained anew using 1.58-bit quantized weights and 8-bit activations. Relative to the initial BitNet version, a number of adjustments have been incorporated, as outlined below.

\paragraph{Quantization Function.} To ensure that the weight values are limited to -1, 0, or +1, we employ an \emph{absmean} quantization technique. This method initially normalizes the weight matrix by dividing it by its mean absolute value and subsequently applies rounding such that each element maps to its closest integer among \{-1, 0, +1\}:

\begin{equation}
    \widetilde{W} = \mathrm{RoundClip}(\frac{W}{\gamma + \epsilon}, -1, 1),
\end{equation}

\begin{equation}
    \mathrm{RoundClip}(x, a, b) = \max(a, \min(b, \mathrm{round}(x))),
\end{equation}

\begin{equation}
    \gamma = \frac{1}{nm}\sum_{ij} |W_{ij}|.
\end{equation}

The method used for activation quantization mirrors that of BitNet, with the key difference that activations are not rescaled to $[0, Q_b]$ prior to non-linear transformation. Instead, all activations are linearly mapped to $[-Q_b, Q_b]$ for each token, thus eliminating the need for zero-point quantization. This approach streamlines both implementation and system optimization, and our empirical analysis shows it results in negligible performance loss.

\paragraph{LLaMA-alike Components.}

The LLaMA architecture~\cite{llama, llama2} has emerged as the standard blueprint for open-source LLMs. To align with trends in the open-source ecosystem, our configuration of \text{BitNet b1.58} leverages LLaMA-inspired design choices. Concretely, the model utilizes RMSNorm~\cite{rmsnorm}, SwiGLU~\cite{swiglu}, and rotary embedding~\cite{rope}, while omitting all bias terms. This compatibility allows \text{BitNet b1.58} to be readily incorporated into leading open-source frameworks (such as Huggingface, vLLM~\cite{vllm}, and llama.cpp\footnote{\url{https://github.com/ggerganov/llama.cpp}}) with minimal modification required.

\section{Results}

We conducted a comparative analysis between \text{BitNet b1.58} and our in-house reproduced FP16 \text{LLaMA LLM} at multiple model scales. To maintain fairness in comparison, both models were pre-trained on the RedPajama dataset~\cite{redpajama} using 100 billion tokens. We assessed zero-shot capabilities across several language understanding benchmarks, namely ARC-Easy~\cite{arc}, ARC-Challenge~\cite{arc}, Hellaswag~\cite{hellaswag}, Winogrande~\cite{winoGrande}, PIQA~\cite{piqa}, OpenbookQA~\cite{openbookqa}, and BoolQ~\cite{boolq}. In addition, we recorded validation perplexity scores on the WikiText2~\cite{wikitext2} and C4~\cite{c4} datasets.

We also analyzed and contrasted both the GPU memory footprint and inference latency for \text{LLaMA LLM} and \text{BitNet b1.58}. All measurements were obtained through the FasterTransformer\footnote{\url{https://github.com/NVIDIA/FasterTransformer}} framework, a high-performance library tailored for large language model inference on GPUs. For \text{BitNet b1.58}, the Ladder~\cite{ladder} 2-bit kernel was incorporated. We focused on reporting the average inference time required per generated token, as it constitutes the principal inference bottleneck.

Table~\ref{tab:ppl} compiles results for perplexity as well as computational and memory resource demands for \text{BitNet b1.58} in comparison to \text{LLaMA LLM}. The findings indicate that, starting from the 3B parameter configuration, \text{BitNet b1.58} achieves perplexity on par with its full precision \text{LLaMA LLM} counterpart, while providing 2.71$\times$ faster inference and reducing GPU memory usage by a factor of 3.55. Notably, the 3.9B variant of \text{BitNet b1.58} achieves 2.4$\times$ speedup and requires only 30\% of the GPU memory, while substantially outperforming the \text{LLaMA LLM} 3B in task performance.

Detailed zero-shot task accuracies are presented in Table~\ref{tab:zero-shot}. Evaluations were conducted following the established \emph{lm-evaluation-harness}\footnote{\url{https://github.com/EleutherAI/lm-evaluation-harness}} methodology. Results demonstrate a diminishing performance difference between \text{BitNet b1.58} and \text{LLaMA LLM} as model sizes scale up. Critically, parity in performance with the full precision baseline is reached by \text{BitNet b1.58} at the 3B scale. As with perplexity, downstream task accuracy measurements highlight that \text{BitNet b1.58} 3.9B outperforms \text{LLaMA LLM} 3B while incurring lower latency and memory costs. These findings confirm that \text{BitNet b1.58} offers a Pareto-optimal advancement compared to leading large language models.

\paragraph{Memory and Latency}

We extended our experiments to larger models, scaling up to 7B, 13B, and 70B parameters, and analyzed their associated costs. As shown in Figure~\ref{fig:latency-memory-energy}, both latency and memory usage display clear patterns: the acceleration effect becomes more pronounced with increased model size. Specifically, the \text{BitNet b1.58} model with 70B parameters achieves a 4.1-fold reduction in inference time relative to the \text{LLaMA LLM} reference model. This advantage arises from the fact that the computational time of \emph{nn.Linear} layers increases with model size. Memory usage follows a similar scaling behavior; since embeddings are stored in full precision, their share of total memory requirements decreases for larger models. Both latency and memory assessments were conducted using a 2-bit kernel, indicating further potential for cost reductions through additional optimizations.

\paragraph{Energy}

We further calculated the energy demand associated with arithmetic operations for both \text{BitNet b1.58} and \text{LLaMA LLM}. Our analysis concentrated primarily on matrix multiplications, which are the most resource-intensive component in LLM computations. Figure~\ref{fig:energy} details the breakdown of these energy costs. The computation in \text{BitNet b1.58} predominantly relies on INT8 additions, whereas \text{LLaMA LLM} utilizes both FP16 additions and FP16 multiplications. Based on the energy models from~\cite{energycost, pokebnn}, \text{BitNet b1.58} demonstrates a 71.4-fold decrease in the energy consumed by arithmetic operations in matrix multiplication on 7nm process technology. We additionally reported end-to-end energy expenditures for processing 512 tokens. The data reveal that \text{BitNet b1.58} delivers greater energy efficiency relative to the FP16 \text{LLaMA LLM} baseline as model size increases. This effect is attributed to the growing share of \emph{nn.Linear} operations with model scale, while other cost components contribute less in larger models.

\paragraph{Throughput}

We assess and contrast the throughput of \text{BitNet b1.58} and the \text{LLaMA LLM} containing 70B parameters across two 80GB A100 GPUs, employing pipeline parallelism~\cite{gpipe} to accommodate \text{LLaMA LLM} 70B on the available hardware. The batch size was progressively increased up to the point where GPU memory was fully utilized, with a fixed sequence length of 512. As presented in Table~\ref{tab:throughput}, \text{BitNet b1.58} 70B supports a maximum batch size that is 11 times greater than that of \text{LLaMA LLM}, resulting in throughput that is 8.9 times higher.

\textbf{\text{BitNet b1.58} introduces a novel scaling paradigm relating model size to performance and inference efficiency}. To illustrate this, the equivalence between 1.58-bit and 16-bit models of various sizes is shown, based on results depicted in Figure~\ref{fig:latency-memory-energy} and \ref{fig:energy}.

\begin{itemize}
    \item The 13B parameter BitNet b1.58 model surpasses a 3B FP16 LLM in terms of latency, memory efficiency, and energy usage.
    \item The 30B BitNet b1.58 model demonstrates superior latency, memory, and energy metrics when compared to a 7B FP16 LLM.
    \item The 70B BitNet b1.58 model is more efficient across latency, memory, and energy consumption than the 13B FP16 LLM.
\end{itemize}

\paragraph{Training with 2T Tokens}

The total number of training tokens is a key variable influencing LLM performance. To explore the scalability of \text{BitNet b1.58} concerning token count, we trained a \text{BitNet b1.58} model using 2T tokens, adhering to the data methodology utilized by StableLM-3B~\cite{StableLM-3B-4E1T}, which stands as a leading open-source 3B parameter model. Evaluation for both models was carried out on a composite benchmark that includes Winogrande~\cite{winoGrande}, PIQA~\cite{piqa}, SciQ~\cite{sciq}, LAMBADA~\cite{lambada}, and ARC-easy~\cite{arc}. Zero-shot accuracy is reported in Table~\ref{tab:2t}; for tasks with both accuracy and normalized accuracy metrics, their average is taken. StableLM 3b's results at 2T tokens are drawn directly from its official technical documentation. The outcomes reveal that \text{BitNet b1.58} consistently outperforms across all evaluated tasks, demonstrating that 1.58-bit LLMs retain robust generalization properties.

\section{Discussion and Future Work}

\textbf{1-bit Mixture-of-Experts (MoE) LLMs}

Mixture-of-Experts (MoE) architectures have demonstrated notable efficiency benefits for large language models (LLMs), especially by lowering the number of floating-point operations (FLOPs). Despite these advantages, MoE systems face practical barriers due to substantial memory demands and considerable inter-chip communication requirements, which complicate scaling and deployment. These obstacles can be mitigated through the adoption of 1.58-bit LLMs. With their reduced memory usage, fewer hardware units are necessary to support MoE deployments. In addition, activation transfer between devices is minimized, resulting in much lower communication overhead. Ideally, if model size allows, housing the entire model on a single chip would eliminate these bottlenecks entirely.

\textbf{Native Support for Long Sequence Processing in LLMs}

Efficient handling of long input sequences has emerged as an essential feature for modern LLMs. A significant limiting factor in long sequence processing is the memory overhead from storing key-value (KV) caches. \text{BitNet b1.58} takes an important step toward alleviating this constraint by reducing activation precision from 16 bits to 8 bits, effectively doubling the possible context length under fixed resource budgets. Further, for 1.58-bit LLMs, there is the potential to compress activations even further, potentially down to 4 bits or less, without information loss—this remains a promising area for subsequent research.

\textbf{LLMs on Edge and Mobile Devices}

Deploying LLMs quantized to 1.58 bits presents a valuable opportunity to enable high-performance language modeling on edge and mobile platforms, which are typically restricted by limited memory and compute resources. Through lower memory consumption and reduced energy requirements, 1.58-bit LLMs can operate efficiently in such constrained environments, unlocking a range of applications previously unfeasible due to hardware limits. This advances the scope of edge and mobile devices, empowering them with more sophisticated language model functionalities. Additionally, 1.58-bit LLMs are inherently better suited for execution on CPUs, the primary processing units in these platforms, thereby enhancing the performance and usability of solutions like \text{BitNet b1.58} in mobile and embedded contexts.

\textbf{Specialized Hardware for 1-bit LLMs}

The recent emergence of projects such as~Groq\footnote{\url{https://groq.com/}} has highlighted the advantages and prospects of crafting dedicated hardware architectures, such as LPUs, tailored for LLM workloads. Building on this momentum, we advocate for further initiatives to create novel hardware and system-level designs specifically optimized for the unique requirements of 1-bit LLMs, taking full advantage of the computational efficiencies introduced by BitNet~\cite{bitnet}.